\textbf{结论名称：} 小范围 \(\Rightarrow\) 大范围（子集法判定充要条件）

\textbf{适用条件：} 条件 \(p,q\) 均可用集合描述（如不等式解集、方程解集、点集等），且判断目标为充分、必要或充要条件。

\textbf{变量说明：} 设 \(A=\{x\mid p(x)\}\)，\(B=\{x\mid q(x)\}\)，全集通常根据题目隐含（如 \(\mathbb{R}\)）。\(A\subseteq B\) 表示 \(A\) 是 \(B\) 的子集。

\textbf{核心公式：}
\[
\boxed{A\subseteq B \iff (p\Rightarrow q) \iff p\text{ 是 }q\text{ 的充分条件}}
\]
特别地，若 \(A=B\) 则 \(p\Leftrightarrow q\)（充要）；若 \(A\subsetneqq B\) 则 \(p\) 是 \(q\) 的充分不必要条件；若 \(A\supseteq B\) 则 \(p\) 是 \(q\) 的必要条件。

\textbf{使用场景：} 利用不等式解集之间的包含关系判定充要条件；已知条件的包含关系反求参数取值范围；将逻辑问题转化为集合运算问题。

\textbf{简单例子：} 设 \(p: x^2-3x+2\le 0\)，\(q: x\le 2\)。  
解 \(p\) 得 \(A=[1,2]\)，\(q\) 对应 \(B=(-\infty,2]\)。  
显然 \(A\subseteq B\) 且 \(A\neq B\)，因此 \(p\Rightarrow q\) 为真而 \(q\nRightarrow p\)，故 \(p\) 是 \(q\) 的充分不必要条件。